% prezentace.tex -- Lekce 01: Seznámení s micro:bitem a MicroPythonem
% preamble-prez.tex -- sdílená preamble pro prezentace (beamer)
% Includuje se z ucitel/lekceXX/prezentace.tex jako:
%   % preamble-prez.tex -- sdílená preamble pro prezentace (beamer)
% Includuje se z ucitel/lekceXX/prezentace.tex jako:
%   % preamble-prez.tex -- sdílená preamble pro prezentace (beamer)
% Includuje se z ucitel/lekceXX/prezentace.tex jako:
%   \input{../spolecne/preamble-prez.tex}

\documentclass[aspectratio=169]{beamer}

% --- Jazyk a font ---
\usepackage{polyglossia}
\setmainlanguage{czech}
\usepackage{fontspec}

% --- Téma ---
\usetheme{default}
\usecolortheme{default}

\definecolor{microbit-blue}{HTML}{1E4D8C}
\definecolor{microbit-lightblue}{HTML}{D6E4F7}
\definecolor{code-bg}{HTML}{F0F0F0}

\setbeamercolor{frametitle}{bg=microbit-blue, fg=white}
\setbeamercolor{title}{fg=microbit-blue}
\setbeamercolor{block title}{bg=microbit-blue, fg=white}
\setbeamercolor{block body}{bg=microbit-lightblue}
\setbeamercolor{itemize item}{fg=microbit-blue}
\setbeamercolor{itemize subitem}{fg=microbit-blue!70}
\setbeamerfont{frametitle}{size=\large, series=\bfseries}
\setbeamertemplate{navigation symbols}{}
\setbeamertemplate{footline}{%
  \leavevmode%
  \hbox{%
    \begin{beamercolorbox}[wd=.5\paperwidth,ht=2.5ex,dp=1.125ex,leftskip=6pt]{author in head/foot}%
      \usebeamerfont{author in head/foot}\textcolor{gray}{\small Úvod do Pythonu na Micro:bitu}
    \end{beamercolorbox}%
    \begin{beamercolorbox}[wd=.5\paperwidth,ht=2.5ex,dp=1.125ex,rightskip=6pt,right]{date in head/foot}%
      \usebeamerfont{date in head/foot}\textcolor{gray}{\small\insertframenumber\,/\,\inserttotalframenumber}
    \end{beamercolorbox}}%
}

% --- Česky korektní uvozovky ---
\usepackage{csquotes}
\newcommand{\uv}[1]{\enquote{#1}}

% --- Zdrojový kód (Python) ---
\usepackage{listings}
\lstset{
  language=Python,
  basicstyle=\ttfamily\small,
  backgroundcolor=\color{code-bg},
  frame=single,
  framerule=0pt,
  xleftmargin=6pt,
  xrightmargin=6pt,
  breaklines=true,
  showstringspaces=false,
  keywordstyle=\color{microbit-blue}\bfseries,
  stringstyle=\color{teal},
  commentstyle=\color{gray}\itshape,
  literate=
    {á}{{\'a}}1 {é}{{\'e}}1 {í}{{\'i}}1 {ó}{{\'o}}1 {ú}{{\'u}}1
    {ů}{{\r{u}}}1 {č}{{\v{c}}}1 {š}{{\v{s}}}1 {ž}{{\v{z}}}1
    {Á}{{\'A}}1 {É}{{\'E}}1 {Í}{{\'I}}1 {Ó}{{\'O}}1 {Ú}{{\'U}}1
    {Č}{{\v{C}}}1 {Š}{{\v{S}}}1 {Ž}{{\v{Z}}}1 {ň}{{\v{n}}}1 {ř}{{\v{r}}}1,
}

% --- Zvýraznění pojmů ---
\newcommand{\pojem}[1]{\textbf{\textcolor{microbit-blue}{#1}}}
\newcommand{\pyinline}[1]{\texttt{#1}}

% --- Grafika a diagramy ---
\usepackage{tikz}

% --- Ostatní ---
\usepackage{graphicx}
\usepackage{hyperref}
\hypersetup{colorlinks=true, urlcolor=microbit-blue}


\documentclass[aspectratio=169]{beamer}

% --- Jazyk a font ---
\usepackage{polyglossia}
\setmainlanguage{czech}
\usepackage{fontspec}

% --- Téma ---
\usetheme{default}
\usecolortheme{default}

\definecolor{microbit-blue}{HTML}{1E4D8C}
\definecolor{microbit-lightblue}{HTML}{D6E4F7}
\definecolor{code-bg}{HTML}{F0F0F0}

\setbeamercolor{frametitle}{bg=microbit-blue, fg=white}
\setbeamercolor{title}{fg=microbit-blue}
\setbeamercolor{block title}{bg=microbit-blue, fg=white}
\setbeamercolor{block body}{bg=microbit-lightblue}
\setbeamercolor{itemize item}{fg=microbit-blue}
\setbeamercolor{itemize subitem}{fg=microbit-blue!70}
\setbeamerfont{frametitle}{size=\large, series=\bfseries}
\setbeamertemplate{navigation symbols}{}
\setbeamertemplate{footline}{%
  \leavevmode%
  \hbox{%
    \begin{beamercolorbox}[wd=.5\paperwidth,ht=2.5ex,dp=1.125ex,leftskip=6pt]{author in head/foot}%
      \usebeamerfont{author in head/foot}\textcolor{gray}{\small Úvod do Pythonu na Micro:bitu}
    \end{beamercolorbox}%
    \begin{beamercolorbox}[wd=.5\paperwidth,ht=2.5ex,dp=1.125ex,rightskip=6pt,right]{date in head/foot}%
      \usebeamerfont{date in head/foot}\textcolor{gray}{\small\insertframenumber\,/\,\inserttotalframenumber}
    \end{beamercolorbox}}%
}

% --- Česky korektní uvozovky ---
\usepackage{csquotes}
\newcommand{\uv}[1]{\enquote{#1}}

% --- Zdrojový kód (Python) ---
\usepackage{listings}
\lstset{
  language=Python,
  basicstyle=\ttfamily\small,
  backgroundcolor=\color{code-bg},
  frame=single,
  framerule=0pt,
  xleftmargin=6pt,
  xrightmargin=6pt,
  breaklines=true,
  showstringspaces=false,
  keywordstyle=\color{microbit-blue}\bfseries,
  stringstyle=\color{teal},
  commentstyle=\color{gray}\itshape,
  literate=
    {á}{{\'a}}1 {é}{{\'e}}1 {í}{{\'i}}1 {ó}{{\'o}}1 {ú}{{\'u}}1
    {ů}{{\r{u}}}1 {č}{{\v{c}}}1 {š}{{\v{s}}}1 {ž}{{\v{z}}}1
    {Á}{{\'A}}1 {É}{{\'E}}1 {Í}{{\'I}}1 {Ó}{{\'O}}1 {Ú}{{\'U}}1
    {Č}{{\v{C}}}1 {Š}{{\v{S}}}1 {Ž}{{\v{Z}}}1 {ň}{{\v{n}}}1 {ř}{{\v{r}}}1,
}

% --- Zvýraznění pojmů ---
\newcommand{\pojem}[1]{\textbf{\textcolor{microbit-blue}{#1}}}
\newcommand{\pyinline}[1]{\texttt{#1}}

% --- Grafika a diagramy ---
\usepackage{tikz}

% --- Ostatní ---
\usepackage{graphicx}
\usepackage{hyperref}
\hypersetup{colorlinks=true, urlcolor=microbit-blue}


\documentclass[aspectratio=169]{beamer}

% --- Jazyk a font ---
\usepackage{polyglossia}
\setmainlanguage{czech}
\usepackage{fontspec}

% --- Téma ---
\usetheme{default}
\usecolortheme{default}

\definecolor{microbit-blue}{HTML}{1E4D8C}
\definecolor{microbit-lightblue}{HTML}{D6E4F7}
\definecolor{code-bg}{HTML}{F0F0F0}

\setbeamercolor{frametitle}{bg=microbit-blue, fg=white}
\setbeamercolor{title}{fg=microbit-blue}
\setbeamercolor{block title}{bg=microbit-blue, fg=white}
\setbeamercolor{block body}{bg=microbit-lightblue}
\setbeamercolor{itemize item}{fg=microbit-blue}
\setbeamercolor{itemize subitem}{fg=microbit-blue!70}
\setbeamerfont{frametitle}{size=\large, series=\bfseries}
\setbeamertemplate{navigation symbols}{}
\setbeamertemplate{footline}{%
  \leavevmode%
  \hbox{%
    \begin{beamercolorbox}[wd=.5\paperwidth,ht=2.5ex,dp=1.125ex,leftskip=6pt]{author in head/foot}%
      \usebeamerfont{author in head/foot}\textcolor{gray}{\small Úvod do Pythonu na Micro:bitu}
    \end{beamercolorbox}%
    \begin{beamercolorbox}[wd=.5\paperwidth,ht=2.5ex,dp=1.125ex,rightskip=6pt,right]{date in head/foot}%
      \usebeamerfont{date in head/foot}\textcolor{gray}{\small\insertframenumber\,/\,\inserttotalframenumber}
    \end{beamercolorbox}}%
}

% --- Česky korektní uvozovky ---
\usepackage{csquotes}
\newcommand{\uv}[1]{\enquote{#1}}

% --- Zdrojový kód (Python) ---
\usepackage{listings}
\lstset{
  language=Python,
  basicstyle=\ttfamily\small,
  backgroundcolor=\color{code-bg},
  frame=single,
  framerule=0pt,
  xleftmargin=6pt,
  xrightmargin=6pt,
  breaklines=true,
  showstringspaces=false,
  keywordstyle=\color{microbit-blue}\bfseries,
  stringstyle=\color{teal},
  commentstyle=\color{gray}\itshape,
  literate=
    {á}{{\'a}}1 {é}{{\'e}}1 {í}{{\'i}}1 {ó}{{\'o}}1 {ú}{{\'u}}1
    {ů}{{\r{u}}}1 {č}{{\v{c}}}1 {š}{{\v{s}}}1 {ž}{{\v{z}}}1
    {Á}{{\'A}}1 {É}{{\'E}}1 {Í}{{\'I}}1 {Ó}{{\'O}}1 {Ú}{{\'U}}1
    {Č}{{\v{C}}}1 {Š}{{\v{S}}}1 {Ž}{{\v{Z}}}1 {ň}{{\v{n}}}1 {ř}{{\v{r}}}1,
}

% --- Zvýraznění pojmů ---
\newcommand{\pojem}[1]{\textbf{\textcolor{microbit-blue}{#1}}}
\newcommand{\pyinline}[1]{\texttt{#1}}

% --- Grafika a diagramy ---
\usepackage{tikz}

% --- Ostatní ---
\usepackage{graphicx}
\usepackage{hyperref}
\hypersetup{colorlinks=true, urlcolor=microbit-blue}


\title{Lekce 01: Seznámení s micro:bitem}
\subtitle{Úvod do Pythonu na Micro:bitu}
\date{}

\begin{document}

\begin{frame}
  \titlepage
\end{frame}

% ------------------------------------------------------------------
\begin{frame}{Co je micro:bit?}
  \begin{columns}
    \begin{column}{0.6\textwidth}
      Malý programovatelný počítač navržený pro výuku:
      \begin{itemize}
        \item \textbf{Displej:} 5 × 5 LED diod
        \item \textbf{Tlačítka:} A a B
        \item \textbf{Senzory:} teplota, pohyb (akcelerometr), magnetismus
        \item \textbf{Zvuk:} reproduktor a mikrofon (micro:bit v2)
        \item \textbf{Komunikace:} Bluetooth, radio, USB
      \end{itemize}
      \medskip
      Programuje se v \textbf{MicroPythonu} — zjednodušené verzi Pythonu.
    \end{column}
    \begin{column}{0.35\textwidth}
      \begin{center}
        \begin{tikzpicture}
          % Jednoduchá schematická kresba micro:bitu
          \draw[rounded corners=4pt, thick, fill=microbit-lightblue]
            (0,0) rectangle (3,4);
          % LED matrix
          \foreach \x in {0.4,0.85,1.3,1.75,2.2}
            \foreach \y in {2.2,2.6,3.0,3.4,3.8}
              \fill[microbit-blue] (\x,\y) circle (0.12);
          % Tlačítka
          \draw[fill=gray!30, rounded corners=2pt] (0.1,1.2) rectangle (0.7,1.6);
          \node at (0.4,1.4) {\tiny A};
          \draw[fill=gray!30, rounded corners=2pt] (2.3,1.2) rectangle (2.9,1.6);
          \node at (2.6,1.4) {\tiny B};
          % USB
          \draw[fill=gray!50] (1.1,0) rectangle (1.9,0.3);
          \node[font=\tiny] at (1.5,0.15) {USB};
        \end{tikzpicture}
      \end{center}
    \end{column}
  \end{columns}
\end{frame}

% ------------------------------------------------------------------
\begin{frame}{Vývojové prostředí: Thonny}
  \pojem{Thonny} je program, ve kterém budeme psát a spouštět kód.

  \bigskip
  \begin{columns}
    \begin{column}{0.48\textwidth}
      \begin{block}{Horní okno — editor}
        Zde píšeme \textbf{program} (skript).\\
        Spuštění: tlačítko \textbf{Run} nebo \textbf{F5}.
      \end{block}
    \end{column}
    \begin{column}{0.48\textwidth}
      \begin{block}{Dolní okno — Shell (REPL)}
        Zde vidíme výstupy programu\\
        a můžeme psát příkazy \uv{za jízdy}.
      \end{block}
    \end{column}
  \end{columns}

  \bigskip
  \textbf{Nastavení interpretu:} Nástroje → Nastavení → Interpreter
  → \uv{MicroPython (BBC micro:bit)}
\end{frame}

% ------------------------------------------------------------------
\begin{frame}{Firmware — proč a jak}
  \pojem{Firmware} je software přímo v micro:bitu — bez něj micro:bit pythonu nerozumí.

  \bigskip
  Micro:bity \textbf{sdílí více skupin} → firmware může chybět nebo být jiný.

  \medskip
  \begin{block}{Postup (jednou za začátku hodiny)}
    \begin{enumerate}
      \item Připoj micro:bit přes USB
      \item Nastav interpreter v Thonny
      \item Pokud Thonny nabídne instalaci → \textbf{klikni Instalovat}
      \item Ověř: v dolním panelu se zobrazí \texttt{MicroPython v...}
    \end{enumerate}
  \end{block}

  \medskip
  \textcolor{gray}{\small Instalace trvá asi 10 sekund. Není to chyba — je to normální postup.}
\end{frame}

% ------------------------------------------------------------------
\begin{frame}{Co je program?}
  \begin{columns}
    \begin{column}{0.55\textwidth}
      \pojem{Program} (nebo \pojem{skript}) je posloupnost příkazů,
      které počítač vykoná jeden po druhém — odshora dolů.

      \bigskip
      Počítač nedělá nic \uv{navíc} — dělá přesně to, co napíšeme.

      \bigskip
      \textcolor{gray}{\small Skript = program uložený v textovém souboru.
      V Pythonu mají soubory příponu \texttt{.py}}
    \end{column}
    \begin{column}{0.4\textwidth}
      \begin{tikzpicture}
        \foreach \i/\cmd in {
          0/{příkaz 1},
          1/{příkaz 2},
          2/{příkaz 3}
        }{
          \draw[fill=microbit-lightblue, rounded corners=3pt]
            (0, -\i*1.1) rectangle (3.2, -\i*1.1+0.8);
          \node[font=\small] at (1.6, -\i*1.1+0.4) {\cmd};
          \ifnum\i<2
            \draw[->, thick, microbit-blue] (1.6,-\i*1.1) -- (1.6,-\i*1.1-0.25);
          \fi
        }
      \end{tikzpicture}
    \end{column}
  \end{columns}
\end{frame}

% ------------------------------------------------------------------
\begin{frame}[fragile]{Funkce a její volání}
  \pojem{Funkce} je pojmenovaná akce — dáme jí jméno a zavoláme ji.

  \bigskip
  \begin{lstlisting}
display.scroll("Ahoj!")
  \end{lstlisting}

  \bigskip
  \begin{itemize}
    \item \pyinline{display.scroll} — \pojem{jméno funkce}
          (\uv{na displeji: rolovej})
    \item \pyinline{( )} — závorky říkají \uv{tuto funkci teď zavolej}
    \item \pyinline{"Ahoj!"} — \pojem{argument}: co funkci předáme
  \end{itemize}

  \bigskip
  \textcolor{gray}{\small Analogie: funkce je jako tlačítko na mikrovlnce.
  Argument je čas, který nastavíme.}
\end{frame}

% ------------------------------------------------------------------
\begin{frame}[fragile]{Náš první program}

  \begin{lstlisting}
from microbit import *

display.scroll("Ahoj!")
  \end{lstlisting}

  \bigskip
  \begin{itemize}
    \item \pyinline{from microbit import *} — načti všechny funkce pro micro:bit
          (tímto začíná každý program)
    \item \pyinline{display.scroll("Ahoj!")} — rolovej text přes displej
  \end{itemize}

  \pause
  \bigskip
  \begin{block}{Vyzkoušej}
    Co se stane, když použiješ \pyinline{display.show()} místo \pyinline{display.scroll()}?
  \end{block}
\end{frame}

% ------------------------------------------------------------------
\begin{frame}[fragile]{\texttt{show} vs.\ \texttt{scroll}}

  \begin{columns}
    \begin{column}{0.48\textwidth}
      \begin{block}{\pyinline{display.show("A")}}
        Zobrazí \textbf{jeden znak} (nebo obrázek) najednou.\\[4pt]
        Rychlé, statické.
      \end{block}
    \end{column}
    \begin{column}{0.48\textwidth}
      \begin{block}{\pyinline{display.scroll("Ahoj!")}}
        Roluje text přes displej\\
        znak po znaku.\\[4pt]
        Pomalejší, pro delší texty.
      \end{block}
    \end{column}
  \end{columns}

  \bigskip

  \begin{lstlisting}
from microbit import *

display.show(Image.HEART)   # ukaze obrazek srdce
  \end{lstlisting}

  \medskip
  Micro:bit má předdefinované obrázky: \pyinline{Image.HEART}, \pyinline{Image.HAPPY},
  \pyinline{Image.SAD}, \pyinline{Image.SURPRISED} \ldots
\end{frame}

% ------------------------------------------------------------------
\begin{frame}{Shrnutí — co jsme se naučili}
  \begin{itemize}
    \item[\checkmark] Micro:bit = malý počítač s LED displejem, tlačítky a senzory
    \item[\checkmark] Thonny = prostředí pro psaní a spouštění programů
    \item[\checkmark] Firmware = software v micro:bitu; sdílené micro:bity ho mohou potřebovat nainstalovat
    \item[\checkmark] Program = posloupnost příkazů vykonaných odshora dolů
    \item[\checkmark] Funkce = pojmenovaná akce; voláme ji závorkami
    \item[\checkmark] Argument = hodnota, kterou funkci předáváme
    \item[\checkmark] \pyinline{from microbit import *} — začíná každý program
    \item[\checkmark] \pyinline{display.scroll()}, \pyinline{display.show()} — zobrazení na displeji
  \end{itemize}
\end{frame}

\end{document}
